\documentclass[a4paper,12pt]{article}
\usepackage[spanish]{babel}
\usepackage[utf8]{inputenc}
\usepackage{geometry}
\usepackage{titling}
\geometry{margin=2.5cm}

\pretitle{}
\posttitle{}

\title{
\vspace{-2cm}
\begin{center}
\LARGE
\rule{\textwidth}{0.5pt} \\[0.3em]
\bfseries
Introducción a las Metodologías Ágiles\\[0.2em]
\large Enfoque en Scrum y Lean Six Sigma
\\[0.3em]
\rule{\textwidth}{0.5pt}
\end{center}
\vspace{-1.5cm}
}

\author{}
\date{}

\begin{document}

\maketitle

\section{Introducción}
Las metodologías ágiles son enfoques de trabajo que buscan mejorar la eficiencia, la adaptabilidad y la calidad en el desarrollo de proyectos. Se caracterizan por promover la colaboración, la entrega continua de resultados y la capacidad de responder rápidamente a cambios.

Dentro de este conjunto, dos metodologías ampliamente utilizadas son \textbf{Scrum} y \textbf{Lean Six Sigma}, cada una con objetivos y enfoques particulares.

\section{Scrum}
Scrum es una metodología ágil orientada principalmente al desarrollo de productos complejos, especialmente en el ámbito del software. Se basa en ciclos de trabajo cortos llamados \textit{sprints}, que suelen durar entre 1 y 4 semanas.

\subsection{Características principales}
\begin{itemize}
    \item Trabajo iterativo e incremental.
    \item Entregas frecuentes de valor.
    \item Fuerte enfoque en el trabajo en equipo.
\end{itemize}

\subsection{Roles en Scrum}
\begin{itemize}
    \item \textbf{Product Owner}: define los requisitos y prioridades.
    \item \textbf{Scrum Master}: facilita el proceso y elimina obstáculos.
    \item \textbf{Equipo de desarrollo}: realiza el trabajo técnico.
\end{itemize}

\subsection{Eventos principales}
\begin{itemize}
    \item \textbf{Sprint}: ciclo de trabajo.
    \item \textbf{Daily Scrum}: reunión diaria de seguimiento.
    \item \textbf{Sprint Review}: revisión de lo realizado.
    \item \textbf{Sprint Retrospective}: análisis para mejorar el proceso.
\end{itemize}

\section{Lean Six Sigma}
Lean Six Sigma es una metodología que combina dos enfoques:
\begin{itemize}
    \item \textbf{Lean}: eliminación de desperdicios.
    \item \textbf{Six Sigma}: reducción de la variabilidad y mejora de la calidad.
\end{itemize}

Su objetivo principal es optimizar procesos, reducir errores y aumentar la eficiencia.

\subsection{Principios básicos}
\begin{itemize}
    \item Enfoque en el cliente.
    \item Mejora continua.
    \item Toma de decisiones basada en datos.
\end{itemize}

\subsection{Ciclo DMAIC}
Lean Six Sigma utiliza un proceso estructurado llamado DMAIC:
\begin{itemize}
    \item \textbf{Definir}: identificar el problema.
    \item \textbf{Medir}: recopilar datos.
    \item \textbf{Analizar}: encontrar causas.
    \item \textbf{Mejorar}: implementar soluciones.
    \item \textbf{Controlar}: mantener las mejoras.
\end{itemize}

\section{Comparación general}
\begin{itemize}
    \item Scrum se enfoca en la gestión de proyectos y el desarrollo ágil.
    \item Lean Six Sigma se centra en la optimización de procesos y la calidad.
    \item Scrum es más flexible y adaptativo.
    \item Lean Six Sigma es más estructurado y basado en datos.
\end{itemize}

\section{Conclusión}
Ambas metodologías buscan mejorar el rendimiento y la eficiencia, pero lo hacen desde enfoques distintos. Scrum es ideal para entornos dinámicos donde los requisitos cambian con frecuencia, mientras que Lean Six Sigma es más adecuado para procesos que requieren control, análisis y mejora continua.

\end{document}
